\documentclass[10pt]{article}

\usepackage[utf8]{inputenc}

\usepackage[T1]{fontenc}

\usepackage{amsmath}

\usepackage{amsfonts}

\usepackage{amssymb}

\usepackage[version=4]{mhchem}

\usepackage{stmaryrd}



\begin{document}

«ХИ-КВАДРАТ» РАСПРЕДЕЛЕНИЕ, $\chi^{2}-p$ а с П p еЛ е л е ни е, - непрерывное, сосредоточенное на положительной полуоси $(0, \infty)$ распределение вероятностей с плотностью



\[

p(x)=\frac{1}{2^{\frac{n}{2}} \Gamma\left(\frac{n}{2}\right)} e^{-\frac{x}{2} x^{\frac{n}{2}-1}},

\]



где $\Gamma(\alpha)$ - гамма-функция, а положительный целочисленный параметр $n$ наз. числом степеней свободы. «X.-к.» р. представляет собой частный случай ааммараспределения и обладает всеми свойствами последнего. Функция распределения «X.-к.» р. есть неполная гаммафункция, характеристич. функция выражается формулой



\[

\varphi(t)=(1-2 i t)^{-n / 2},

\]



математич. ожидание и дисперсия равны, соответственно, $n$ и $2 n$. Семейство «X.-к.» р. замкнуто относительно операции свертки.



«X.-к.» р. с $n$ степенями свободы может быть выведено как распределение суммы $\chi_{n}^{2}=X_{1}^{2}+\ldots+X_{n}^{2}$ квадратов независимых случайных величин $X_{1}, \ldots, X_{n}$, имеющих одинаковое нормальное распределение с нулевым математич. ожиданием и единичной дисперсией. Эта связь с нормальным распределением определяет ту роль, к-рую «Х.-к.» р. играет в теории вероятностей И математич. статистике.



Многие распределения определяются посредством «X.-к.» р. Таковы, напр., распределение случайной величины $\sqrt{\chi_{n}^{2}}$ - длины случайного вектора $\left(X_{1}, \ldots, X_{n}\right)$ с независимыми, нормально распределенными компонентами (иногда наз. «х и»р а с п р е д е л е н и е м, см. также частные случаи - Максвелла распределение, Рәлея распределение); Стьюдента распределение, Фишера $F$-pacnределение. В математич. статистике эти распределения вместе с «X.-к.» р. описывают выборочные распределения различных статистик от нормально распределенных результатов наблюдений и используются для построения интервальных статистич. оценок и статистических критериев. Особую известность в связи с «.-к.» р. получил «xи-квадрат» критерий, основанный на т. н. «хи-квадрат» статистике Пирсона.



Имеются подробные удобные для статистич. расчетов таблицы «X.-к.» р. При больших $n$ используют анпроксимации посредством нормального распределения: нашр., согласно центральвой предельной теореме распределение нормированной велитины $\frac{x_{n}^{2}-n}{\sqrt{2 n}}$ сходится к стандартному нормальному распределению. Более точна аппроксимация



\[

\begin{array}{r}

\mathrm{P}\left\{\chi_{n}^{2}<x\right\} \rightarrow \Phi(\sqrt{2 x})-\Phi(\sqrt{2 n-1}) \\

\text { при } n \rightarrow \infty,

\end{array}

\]



где $\Phi(x)$ - функция стандартного нормального распределения. ure.



См. также Нецентральное «хи-квадрат» распределе-



Јіт.: [1] К р а м е р Г., Математические методы статистики, пер. с англ., 2 изд., М., $1975 ;$ [2] К е н д а л л М. Дж., С т ь ю а р т А., Теория распределений, пер. с англ., М., 1966; [a. o], 1969; [3] Бо ль т е в J. H., С м и Таблицы математической статистики, 3 изд., М., 1983 .



ХИЛЛА УРАВНЕНИЕ - обыкновенное дирохоров. циальное уравнение 2-го порядка



\[

w^{\prime \prime}(z)+p(z) w(z)=0

\]



c периодич. Функцией $p(z)$; все величины могут быть комплексными. Уравнение названо по имени Дж. Хил- ла [1], к-рый, изучая движение Луны, получил уравнение



\[

w^{\prime \prime}(z)+\left(\theta_{0}+2 \sum_{r=1}^{\infty} \theta_{2 r} \cos 2 r z\right) w(z)=0

\]



с действительными числами $\theta_{0}, \theta_{2}, \theta_{4}, \ldots$, причем ряд $\sum_{r=1}^{\infty}\left|\theta_{2 r}\right|$ сходится.



Дж. Хилл дал метод решения Х. у. с использованием определителей бесконечного порядка. Это явилось толчком для создания теории таких определителей и, далее, для создания Э. Фредгольмом (E. Fredholm) теории интегральных уравнений. Для $\mathrm{X}$. у. ставятся прежде всего задачи устойчивости решений, существования или отсутствия периодич. решений. Если в действительном случае в Х. у. ввести параметр:



\[

x^{\prime \prime}+\lambda p(t) x=0 \text {, }

\]



то, как установил А. М. Ляпунов [2], существует такая бесконечная последовательность



\[

\begin{gathered}

\ldots<\lambda_{-1} \leqslant \lambda_{0}=0<\lambda_{1} \leqslant \lambda_{2}<\ldots<\lambda_{2 n-1} \leqslant \\

\leqslant \lambda_{2 n}<\lambda_{2 n+1} \leqslant \ldots,

\end{gathered}

\]



что при $\lambda \in\left(\lambda_{2 n}, \lambda_{2 n+1}\right) \mathrm{X}$. у. устойчиво, а при $\lambda \in\left[\lambda_{2 n-1}, \lambda_{2 n}\right] \mathrm{X} . \mathrm{y}$. неустойчиво. При әтом $\lambda_{4 n}$ и $\lambda_{4 n+3}$ являются собственными значениями периодич. краевой задачи, а $\lambda_{4 n+1}$ и $\lambda_{4 n+2}-$ собственными значениями полупериодич. краевой задачи. Хорошо изучена теория Х.У. (см. [3]).



Jum.: [1] H i l l G., "Acta math.», 1886, v. 8, p. 1; [2] Л я п yн о в А. М., Собр. соч., т. 2, М., 1956, с. $407-09$; [3] Я к у б ов и ч В. А., С т а р жк и н с к и й В. М., Јинейные дифференциальные уравнения с периодическими коэфффииентами и их ХИН ЧИНА ИНТЕГРАЛ - обобщение узкого Данжуа интеграла, введенное А. Я. Хинчиным [1].



Функция $f(x)$ наз. и н т е г р и р у е м о й в с м ы сл е X и н ч и н а на $[a, b]$, если она интегрируема широким интегралом Данжуа и ее неопределенный интеграл почти всюду дифференцируем. Иногда интеграл Хинчина наз. ин т е г р а лом Д ан жу а- X инч И н а.



Лum.: [1] Х ин ч ин А. Я., “C. r. Acad. sci.", 1916, t. 162, p. 287-91; [2] Х и н ч и н А.'Я., «Матем. сб.», 1918, т. 30, М., 1966 , с. $171-73 ;$ [4] С a к с С., Теория интеграла, пер. с англ., М., 1949 , с. 370 . T. П. Лукашенко.



ХИНЧИНА НЕРАВЕНСТВО д л я н $е$ з а в Я с и мы ф ун к ц и й - оценка в $L_{p}$ суммы независимых функций. Пусть $f_{k}$-система независимых функций и для нек-рого $p>2$



\[

\sup _{k}\left\|f_{k}\right\|_{L_{p}}<\infty, \quad \int_{0}^{1} f_{k}(t) d t=0 .

\]



Тогда



Если



\[

\left\|\sum_{k=0}^{\infty} c_{k} f_{k}\right\|_{L_{p}} \leqslant M\left(\sum_{k=1}^{\infty} c_{k}^{2}\right)^{1 / 2}

\]



\[

\sum_{k=1}^{\infty} c_{k}^{2}<\infty, \text { a } r_{k}=\operatorname{sign} \sin 2^{k} \pi t

\]



\begin{itemize}

  \item функции Радемахера и

\end{itemize}



\[

f(t)=\sum_{k=1}^{\infty} c_{k} r_{k}(t)

\]



то для любого $p>0$



\[

\begin{gathered}

A_{p}\left(\sum_{k=1}^{\infty} c_{k}^{2}\right)^{1 / 2} \leqslant\left(\int_{0}^{1}|f(t)| p d t\right)^{1 / p} \leqslant \\

\leqslant B_{p}\left(\sum_{k=1}^{\infty} c_{k}^{2}\right)^{1 / 2},

\end{gathered}

\]



где $B_{p}=O(\sqrt{p})$ при $p \rightarrow \infty$. Это неравенство было установлено А. Я. Хинчиным [1]. Точное значение $A_{1}$ равно $1 / 2$.





\end{document}